\section{Introduction}

Modern online games produce telemetry data that records player movement, spatial conditions, interactions, and in-game events in temporal order. The global game market is forecast to reach approximately USD 188.8 billion with about 3.58 billion players in 2025~\cite{newzoo2025globalgames}, and telemetry at this scale is an important data resource for a wide range of analyses. This data is more than post-match aggregate statistics: it also captures the spatial and temporal context in which play unfolds. Battle royale games are an especially clear case of that structure. Players or teams enter a large battlefield, secure resources, move through a shrinking activity space, and compete to survive until the end.

We focus on \textit{PUBG: Battlegrounds}, a leading and representative battle royale title, as the case study for this work. A single PUBG match places up to 100 players on a large map, and the area in which they can safely remain shrinks over the course of the match~\cite{pubgGameInfo}.

KRAFTON, the developer, publishes official match telemetry that records phase changes alongside player positions, damage, knockdowns, kills, and eliminations in temporal order~\cite{pubgTelemetryEvents}. This stream is spatio-temporal by nature: teams are moving objects, the safe and blue zones are dynamic regions of the map, and combat and elimination are events anchored to a particular time and place. What makes it both special and difficult is that a team's survival risk does not hold steady within a match; it rises and falls from one phase to the next. Post-match aggregate statistics such as total kills, total damage, or final travel distance flatten this away, so they cannot tell us when a team's risk rose or which spatial condition was behind it.

Most public benchmarks and prediction studies on PUBG take a different angle: they predict a match's final placement or final rank from post-match aggregate statistics~\cite{kagglePubgFinishPlacement,sen2021pubgFinalRank,lee2025pubgRank}. These methods suit final-outcome prediction, but by compressing the staged survival process into a single label they cannot say when a team became vulnerable, which phase raised the spatial pressure on it, or how nearby-team configurations and the moving safe zone changed its odds of elimination. Recent phase-based PUBG survival work moves closer to the problem we study; in the Related Work section, we use that work as the direct point of departure. We take a different route at the representation level. We reconstruct PUBG telemetry as a spatio-temporal database with dynamic inter-team relations, and we cast survival as a phase-level \textit{next-phase elimination hazard}, framed as a discrete-time event-history problem~\cite{allison1982discrete,singer1993discrete}. We keep this contrast brief here and return to it in the dedicated Related Work section.

Three research questions organize the study:
\begin{description}
  \item[\textbf{RQ1.}] How can PUBG telemetry be represented as a spatio-temporal database that captures moving teams, dynamic safe and blue zones, spatial events, and the phase structure of a match?
  \item[\textbf{RQ2.}] As the safe zone contracts, how do team trajectories and survival conditions change from one phase to the next?
  \item[\textbf{RQ3.}] Do spatio-temporal features and dynamic inter-team graph relations contribute to predicting next-phase team elimination risk, beyond strong tabular baselines?
\end{description}
